% Nouveau EGA V (cote 25): archive cote 25, site page 28 (PDF page 29).
% typescript. Draft compared: none.
% AI-transcribed by Claude Opus 5.5 (claude-opus-5-5, adaptive thinking, effort high), Message Batches API, batch ega25b (msgbatch_01A7n5YuXubUPNK1zsepi8RG), 2026-11-13; not yet checked by eye.
\scanpage{25/28}
\mspage{29}
pour tout t \hand{$\in$}U,  sa fibre \inserted{$Y_{\xi}$(1)} en  t satisfasse la propriété P.  Lorsqu'on ne
dispose pas à priori de \hand{$\xi$}$_{o}$ , \xout{xxx} mais que k(s) est infini, on pourra combi-%
ner les deux remarques précédentes pour obtenir une section hyperplane
sur un voisinage ouvert de s, ayant la propriété précédente. \xout{Dixxxxx} Enfin,
utilisant \deleted{4.8.}\inserted{le N$^{o}$5} , on a un critère permettant d'affirmer que (X resp. F étant
supposé plat sur S, ce qui permet d'appliquer \deleted{4.7}\inserted{loc. cit.}.)  Y\hand{$_{\xi}$} \inserted[typed]{resp. G\hand{$_{\xi}$}} est également
plat sur S. On peut alors, remplaçant X par Y\hand{$_{\xi}$} , itèrer la construction
précédente, ce qui permett\xout{xxxxxx} \inserted[typed]{par exemple\hand{,} sous certaines conditions} de construire de proche en proche une
``multisection'' \inserted[typed]{S'} de X au dessus d'un voisinage ouvert U d'un point donné s,
telle\xout{x} que S'\hand{$\longrightarrow$}U soit fini, plat et à fibres satisfaisant la propriété P.
Lorsque k(s) est fini, on peut être contraint de faire un changement de bas%
% NOTE: right edge of the page cut; last letter of "base" not visible
étale surjectif S'\hand{$\longrightarrow$} U (U voisinage ouvert de s) avant de pouvoir appliquer
les constructions précédentes; en effet, sous les conditions du début de
6.1. , si k est fini, il n'existe pas nécessairement de point rationnel sur
k dans  l'ouvert non vide U, mais \xout{l'ensemble des points fermés de U est}
il existe certainement un point fermé de U, donc un point à valeurs dans
une extension f inie k' (nécessairement séparable) de k; lorsque k = k(s),
on peut donc, quitte à faire une extension \inserted[typed]{finie} étale convenable \inserted[typed]{S'} sur un voisinag%
% NOTE: right edge cut; end of "voisinage" not visible
ouvert U de  s, correspondant à l'extension résiduelle k'/k\xout{xxxxxxxxx}
i.e. telle que  S$'_{s}$ \hand{$\xrightarrow{\sim}$}\inserted{\unsure{k(s)-isomorphisme}}Spec(\over{k}{k'}) , se ramener à la situation favorable \xout{par}
en l'unique point s'\hand{$\in$}S' au dessus de s, par le changement de base S'\hand{$\to$} S.

Donnons quelques exemples d'énoncés obtenus par ce\over{s}{tte} méthode\deleted{s}. \deleted{Il y a}
\deleted{lieu cependant, de faire une remarque préliminaire, pour l'élimination dan}%
% NOTE: right edge cut ("dans")
\deleted{\uline{\xout{Propositionxxxxxxxxxx}}}
\deleted{ces constructions de l'hypothèse faite dans 4.6. \deleted{et} 4.7. \inserted{et 4.8.} que  aucune compo}%
\deleted{sante irréductible géométrique d\xout{x}'un supp F$_{s}$ n'est contenue\hand{\unsure{,}}dans une sectio}%
\deleted{hyperplane (definie sur $\overline{\text{k(s)}}$ ): on peut (et doit) en effet énoncer 4.6.\hand{\unsure{/}}}
\deleted{et 4.7. sous forme légèrement plus générale, en se donnant une partie ouver}%
\deleted{te W de P\hand{$\check{}$} , \xout{ex}  en remplaçant dans la condition envisagée  ``hyperplan'' par}
\margin{28}
